\section{Passing to three-dimensional space}
\label{sec:geometry}

It remains to remove the finite coordinate in \cref{stack:cyclic}.
We construct a finite tile in $\Z^3$ whose unit-cube thickening also
excludes periodic tilings with arbitrary real translation vectors.
Quotient-to-lattice reductions are studied by \citet[Theorem~1.1]{MSSReduction},
building on earlier rigid tile constructions such as
\citet[Lemma~9.3]{GTTwoTiles}. The need to constrain real translation
vectors also appears in the Euclidean reduction of
\citet[Theorem~2.1 and Lemma~2.2]{GTPeriodic}. The argument below
proves both required statements directly for the same finite tile
and its unit-cube thickening.

Let
\[
 \pi:\Z^3\longrightarrow\Gamma,\qquad
 \pi(a_1,a_2,a_3)=(a_1,a_2,a_3\bmod Q),
 \qquad w=(0,0,Q).
\]
Thus $\ker\pi=\Z w$.  Translate the tile $F$ of \cref{stack:cyclic}
so that $0\in F$, and choose a finite set $U\subset\Z^3$ containing
$0$ on which $\pi$ is a bijection onto $F$.  In particular,
\begin{equation}
 U\cap\ker\pi=\{0\},\qquad w\notin U.
 \label{geom:representatives-U}
\end{equation}
Translating $F$ changes neither tileability nor the existence of a
fully periodic tiling.

\subsection{A residue transversal with many differences}

\begin{lemma}[Rigid representatives]
\label{geom:random-representatives}
For a sufficiently large integer $m\geq2$, one can choose
\[
 T_0=\{t(v):v\in\{0,\ldots,m-1\}^3\},\qquad
 t(v)\equiv v\pmod m,
\]
such that, with $T_*=T_0\setminus\{t(0)\}$,
\begin{equation}
 T_*-T_*\supset
 \{d\in\Z^3:\norm{d}_\infty\leq m,
                         \ d\notin m\Z^3\}.
 \label{geom:difference-property}
\end{equation}
\end{lemma}

\begin{proof}
Choose independently, for each residue $v$,
\[
 t(v)=v+mh_v,\qquad
 h_v\text{ uniform in }\{-2,-1,0,1,2\}^3.
\]
Fix a requested vector $d$ on the right-hand side of
\cref{geom:difference-property}.  Translation by $d\bmod m$ on
$(\Zmod{m})^3$ is a permutation without fixed points.  A cycle of
length $L\geq2$ contains $\lfloor L/2\rfloor\geq L/3$ disjoint
successive pairs.  Selecting these pairs in every cycle gives at
least $m^3/3$ disjoint pairs of residues $(u,v)$ with
$v-u\equiv d\pmod m$.  Discard the pair containing residue $0$, if
there is one.  At least
\[
 k_m=\lfloor m^3/3\rfloor-1
\]
pairs remain, all with both endpoints nonzero.

For such a pair, the equality $t(v)-t(u)=d$ asks for
\[
 h_v-h_u=\frac{d-(v-u)}m.
\]
This is an integer vector whose coordinates have absolute value at
most $2$, since $\norm{d}_\infty\leq m$ and the coordinates of $u,v$
lie between $0$ and $m-1$.  Two independent uniform variables on
$\{-2,-1,0,1,2\}$ have any specified difference of absolute value at
most $2$ with probability at least $3/25$.  Hence each chosen pair
realizes $d$ with probability at least
\[
 c=(3/25)^3=27/15625.
\]
These events are independent because the pairs have disjoint
endpoints.  Failure to realize this $d$ has probability at most
$(1-c)^{k_m}$.  There are at most $(2m+1)^3$ requested vectors, so
the probability that \cref{geom:difference-property} fails is at most
\[
 (2m+1)^3(1-c)^{k_m}.
\]
This tends to zero as $m$ tends to infinity.  Fix $m\geq2$ for which
it is less than one, and then fix a successful choice of the
representatives.  Every witnessing pair avoids residue $0$, as
required.
\end{proof}

Define the second transversal by moving just the marked point:
\[
 T_1=T_*\cup\{t(0)+mw\}.
\]
Both $T_0$ and $T_1$ have exactly one representative of each residue
modulo $m$.  The switch is illustrated schematically in
\cref{geom:fig-switch}.

\begin{figure}[!htbp]
\centering
\begin{tikzpicture}[
  x=1cm,y=1cm,
  every node/.style={font=\small},
  common/.style={draw=black!55,fill=black!16,line width=.5pt},
  marked/.style={draw=BurntOrange,fill=BurntOrange!30,line width=.9pt}
]
  \node[anchor=east] at (0.8,1.35) {$T_0$};
  \node[anchor=east] at (0.8,-0.35) {$T_1$};

  \foreach \y in {1.35,-0.35} {
    \draw[common] (1.55,\y-.16) rectangle (1.87,\y+.16);
    \draw[common] (2.17,\y-.16) rectangle (2.49,\y+.16);
    \node at (2.94,\y) {$\cdots$};
    \draw[common] (3.40,\y-.16) rectangle (3.72,\y+.16);
    \node at (2.64,\y-.51) {$T_*$};
  }

  \draw[marked] (4.84,1.19) rectangle (5.16,1.51);
  \node at (5,0.84) {$t(0)$};

  \draw[draw=black!45,dashed,line width=.6pt]
    (4.84,-.51) rectangle (5.16,-.19);
  \draw[-{Latex[length=2.2mm]},line width=.8pt,BurntOrange]
    (5.40,-.35) -- node[above=3pt,text=black] {$mw$} (9.36,-.35);
  \draw[marked] (9.60,-.51) rectangle (9.92,-.19);
  \node at (9.76,-.86) {$t(0)+mw$};
\end{tikzpicture}
\caption{Schematic of the marked-point switch, without depicting the
actual positions or scale.  The gray marks represent the same common
set $T_*$ in both rows.  The marked point moves from $t(0)$ to
$t(0)+mw$; the dashed outline shows its former location.  The marks
may also denote unit-cube thickenings.  In a grid assembly this switch
gives the kernel-invariance identity \eqref{geom:kernel-invariance}.}
\label{geom:fig-switch}
\end{figure}


Our final discrete tile is
\begin{equation}
 T=T_1\cup\bigcup_{u\in U\setminus\{0\}}(mu+T_0),
 \qquad \Omega=T+[0,1]^3.
 \label{geom:large-tile}
\end{equation}
The finite union defining $T$ is disjoint.  To see this, equality of
two points first forces equality of their residue labels $v$ modulo
$m$.  For $v\ne0$, the points are $mu+t(v)$, and equality forces
the same $u$.  For $v=0$, the unmarked points are
$mu+t(0)$ with $u\in U\setminus\{0\}$, while the marked point is
$mw+t(0)$.  They are all distinct by
\cref{geom:representatives-U}.  In particular, $|T|=|U|m^3$.

Write
\[
 \Omega_i=T_i+[0,1]^3\quad(i=0,1),\qquad
 \Omega_*=T_*+[0,1]^3.
\]
Each $\Omega_i$ has volume $m^3$, and their common subset $\Omega_*$
has positive volume $m^3-1$.  Unit cubes with distinct integer lower
corners have disjoint interiors, so the decomposition of $T$ in
\cref{geom:large-tile} gives the corresponding decomposition of
$\Omega$ up to null boundaries.
We call the translation vector $c$ the center of the copy $c+T_i$
or $c+\Omega_i$.

\subsection{Existence of a tiling}

\begin{proposition}
\label{geom:existence}
The tile $T$ tiles $\Z^3$, and $\Omega$ tiles $\R^3$ up to null sets.
\end{proposition}

\begin{proof}
Let $A\oplus F=\Gamma$ be a tiling from \cref{stack:cyclic}, and put
$B=\pi^{-1}(A)$.  Then $B+w=B$.  Also
$B\oplus U=\Z^3$: for any $g\in\Z^3$, the unique representation of
$\pi(g)$ in $A+F$ determines a unique $u\in U$, and then $g-u\in B$
is the unique possible other summand.

Place copies of $T$ at all translations in $mB$.  Their constituent
copies of $T_0$ or $T_1$ have centers $m(b+u)$, for $b\in B$ and
$u\in U$.  Since $B\oplus U=\Z^3$, there is exactly one such center
at every point of $m\Z^3$.  Its type is $1$ precisely at the centers
$mB$, because the type-$1$ piece corresponds to $u=0$.

The assembly with $T_0$ at every center in $m\Z^3$ tiles $\Z^3$
exactly: each integer vector has a unique residue label $v$ and a
unique expression $ma+t(v)$.  Switching to type $1$ at the centers
$mB$ removes the marked points
\[
 t(0)+mB
\]
and inserts
\[
 t(0)+m(B+w)=t(0)+mB.
\]
The insertion is a bijective relocation, so all multiplicities
remain one.  This proves $mB\oplus T=\Z^3$.  Thickening by unit
cubes gives an almost-everywhere tiling by $\Omega$ with the same
translation set: away from the integer coordinate planes, every
point belongs to the interior of one integer unit cube.
\end{proof}

\subsection{Rigidity for real translation vectors}

We next show why arbitrary real translations cannot produce a
periodic tiling.  The marked-point modification has two roles.
The common part forces the centers of the small pieces onto one
grid, and the moved point then forces invariance under the kernel
direction $w$.

\begin{lemma}[Rounding in the closed box]
\label{geom:rounding}
Let $m\geq2$, and suppose that $\delta\in\R^3$ satisfies
$\norm{\delta}_\infty\leq m$ and $\delta\notin m\Z^3$.  There is
$d\in\Z^3\setminus m\Z^3$ such that
\[
 \norm{d}_\infty\leq m,
 \qquad \norm{\delta-d}_\infty<1.
\]
\end{lemma}

\begin{proof}
If $\delta$ is integral, use $d=\delta$.  Otherwise, keep every
integer coordinate unchanged and round each noninteger coordinate
either down or up.  Both choices lie in $[-m,m]$, since the
endpoints are integers, and both are at distance strictly less than
one from the original coordinate.  In at least one noninteger
coordinate choose a rounding which is not a multiple of $m$.
Such a choice exists because two consecutive integers cannot both
be multiples of $m\geq2$.  The resulting $d$ is not in $m\Z^3$.
Coordinates of $\delta$ equal to $m$ or $-m$ are left unchanged, so
the strict distance estimate also holds on the boundary of the box.
\end{proof}

\begin{lemma}[A periodic assembly has grid centers]
\label{geom:grid}
Suppose copies $c+\Omega_{\iota(c)}$, with types
$\iota(c)\in\{0,1\}$, form an almost-everywhere tiling of $\R^3$.
Assume the assembly, including its types, is invariant under a
full-rank lattice $\Lambda\subset\R^3$.  Then the set $C$ of centers
is $c_0+m\Z^3$ for some $c_0\in\R^3$.
\end{lemma}

\begin{proof}
First, two constituent copies cannot have the same center, because
both contain the positive-volume common part $\Omega_*$.  Thus the
centers and their types are unambiguous.  We claim that distinct
centers $c,c'$ satisfy
\begin{equation}
 \norm{c'-c}_\infty\leq m
 \quad\Longrightarrow\quad c'-c\in m\Z^3.
 \label{geom:nearby-centers}
\end{equation}
If $\delta=c'-c$ violated this assertion, \cref{geom:rounding} would
give a vector $d$ in the difference set in
\eqref{geom:difference-property} with
$\norm{\delta-d}_\infty<1$.  Write $d=t-t'$ for $t,t'\in T_*$.  The
two common-part unit cubes
\[
 c+t+[0,1]^3,\qquad c'+t'+[0,1]^3
\]
have lower corners whose coordinate differences have absolute
value strictly less than one.  Their intersection therefore has
positive volume, contradicting the tiling property.  This proves
\cref{geom:nearby-centers}.

In particular, distinct centers have sup-norm distance at least
$m$: a nonzero vector in $m\Z^3$ cannot have smaller sup norm.
The set $C$ is consequently uniformly separated and locally finite.
The auxiliary closed cubes
\[
 Q_c=c+[-m/2,m/2]^3\qquad(c\in C)
\]
have disjoint interiors.  Each has volume $m^3$, equal to the
volume of either constituent type.

We use periodicity to show that these auxiliary cubes cover.  Let
$P$ be a bounded half-open fundamental parallelepiped for $\Lambda$.
There are finitely many centers modulo $\Lambda$, by local
finiteness; choose representatives $c_1,\ldots,c_r$.  Their types
$\iota_1,\ldots,\iota_r$ are preserved under $\Lambda$.  Integrating
the constituent tiling over $P$ gives
\begin{align*}
 |P|
 &=\int_P\sum_{j=1}^r\sum_{\lambda\in\Lambda}
       \mathbf{1}_{c_j+\lambda+\Omega_{\iota_j}}(x)\,dx\\
 &=\sum_{j=1}^r|\Omega_{\iota_j}|=rm^3.
\end{align*}
For the second equality, translate the integration domains by
$-\lambda$ and use that $\{P-\lambda:\lambda\in\Lambda\}$ partitions
$\R^3$ up to boundaries.  Nonnegative summands justify the exchange
of summation and integration.

The same calculation for the auxiliary cubes gives
\[
 \int_P\sum_{c\in C}\mathbf{1}_{Q_c}(x)\,dx=rm^3=|P|.
\]
The integrand is at most one almost everywhere, because the cubes
have disjoint interiors; their countably many boundaries are null.
It is therefore one almost everywhere on $P$, and periodicity gives
almost-everywhere coverage of $\R^3$.  The family of closed cubes
$\{Q_c\}_{c\in C}$ is locally finite, so its union is closed.  A
point outside the union would have an open neighborhood of positive
volume outside it, contradicting almost-everywhere coverage.  Thus
the cubes cover every point of $\R^3$.

Their intersection graph is connected.  Indeed, if one graph
component were separated from the remaining components, the unions
of their respective cubes would be disjoint nonempty closed sets
covering $\R^3$.  Both unions are closed because they are subfamilies
of a locally finite family of closed sets.  They would disconnect
$\R^3$, which is impossible.  Whenever $Q_c$ and $Q_{c'}$ intersect,
$\norm{c'-c}_\infty\leq m$, so \cref{geom:nearby-centers} puts their
difference in $m\Z^3$.  Following paths in the intersection graph
therefore puts every center in one coset $c_0+m\Z^3$.

No point of that coset can be missing.  The interior of the
side-$m$ cube at a missing grid point is disjoint from all side-$m$
cubes at the other grid points and would be uncovered.  Hence
$C=c_0+m\Z^3$.
\end{proof}

\subsection{Kernel invariance and descent}

\begin{proposition}
\label{geom:no-periodic}
The Euclidean tile $\Omega$ has no almost-everywhere translational
tiling whose translation set is invariant under a full-rank lattice
in $\R^3$.
\end{proposition}

\begin{proof}
Suppose, to the contrary, that $\mathcal A\subset\R^3$ is such a
translation set and that $\mathcal A+\Lambda=\mathcal A$ for a
full-rank lattice $\Lambda$.  Distinct vectors of $\mathcal A$ have
sup-norm distance at least one: otherwise their translates of a fixed
unit cube in $\Omega$ would overlap in positive volume.  Hence
$\mathcal A$ is countable.  Decompose every translated copy of
$\Omega$ using \cref{geom:large-tile}.  A copy at $a\in\mathcal A$
has a type-$1$ constituent centered at $a$ and type-$0$ constituents
centered at $a+mu$ for $u\in U\setminus\{0\}$.  The constituents
form an almost-everywhere tiling and inherit the lattice periods,
including their types.  No two constituents have the same center,
as observed in the proof of \cref{geom:grid}.  Consequently the map
\begin{equation}
 \mathcal A\times U\longrightarrow C,\qquad(a,u)\longmapsto a+mu
 \label{geom:center-decomposition}
\end{equation}
is a bijection, and the type-$1$ centers are exactly $\mathcal A$.

By \cref{geom:grid}, $C=c_0+m\Z^3$.  Translate the entire assembly
by $-c_0$, which does not change its period lattice, and henceforth
write $C=m\Z^3$.  There is a set $B'\subset\Z^3$ such that
$\mathcal A=mB'$.  Let $b(a)=\mathbf{1}_{B'}(a)$; this records the
type at the center $ma$.  The bijection
\eqref{geom:center-decomposition}, divided by $m$, gives
\begin{equation}
 B'\oplus U=\Z^3.
 \label{geom:upstairs-tiling}
\end{equation}

Consider the open unit cell with lower corner $t(0)+ma$, for any
$a\in\Z^3$.  Since all constituent centers and all voxel positions
are integral, a unit cube can meet the interior of this cell only
if it has the same lower corner.  The residues modulo $m$ then
exclude every common-part point $t(v)$ with $v\ne0$.  There are
exactly two possible sources: the unshifted marked point from
type $0$ at $ma$, or the moved marked point from type $1$ at
$m(a-w)$.  The coverage multiplicity on this open cell is therefore
\[
 1-b(a)+b(a-w).
\]
The almost-everywhere tiling makes it one.  Thus $b(a)=b(a-w)$ for
every $a$, or equivalently
\begin{equation}
 B'+w=B'.
 \label{geom:kernel-invariance}
\end{equation}
This also gives invariance under every element of $\ker\pi=\Z w$.

The set $B'$ is fully periodic in the discrete sense.  It is
nonempty, and $\mathcal A=mB'\subset m\Z^3$.  For $a_0\in\mathcal A$
and any $\lambda\in\Lambda$, both $a_0$ and $a_0+\lambda$ belong to
$m\Z^3$, so $\lambda\in m\Z^3$.  Hence
$L=m^{-1}\Lambda$ is a rank-three lattice contained in $\Z^3$ and
preserves $B'$.  Such an integer lattice has finite index: the
matrix formed by an integer basis has nonzero integer determinant,
whose absolute value is its index.  In particular, $L\leq\Z^3$ is
a finite-index period subgroup of $B'$.

We finally descend the exact tiling
\eqref{geom:upstairs-tiling} to the quotient.  Fix $\gamma\in\Gamma$
and any lift $g\in\Z^3$.  For each $u\in U$, there is a
representation
\[
 \gamma=\alpha+\pi(u),\qquad \alpha\in\pi(B'),
\]
if and only if $g-u\in B'$.  One direction follows by applying
$\pi$.  For the other, if $\alpha=\pi(b_0)$ with $b_0\in B'$, then
$g-u-b_0\in\ker\pi$; \cref{geom:kernel-invariance} therefore puts
$g-u$ in $B'$.  By \cref{geom:upstairs-tiling}, exactly one $u$ has
this property, and $\pi$ is injective on $U$.  Thus
\[
 \pi(B')\oplus\pi(U)=\Gamma.
\]
This quotient tiling is fully periodic.  Indeed, $\pi(L)$ preserves
$\pi(B')$, and the surjection
$\Z^3/L\to\Gamma/\pi(L)$ shows that $\pi(L)$ has finite index.
Since $\pi(U)=F$, this contradicts \cref{stack:cyclic}.
\end{proof}

\begin{proof}[Completion of the proof of \cref{thm:main}]
The finite nonempty tile $T$ from \cref{geom:large-tile} and its
closed unit-cube thickening $\Omega$ admit the tilings in
\cref{geom:existence}.  The Euclidean tiling cannot be fully
periodic by \cref{geom:no-periodic}.  If a discrete tiling of $T$
had a finite-index period subgroup in $\Z^3$, thickening its unit
cells would give an almost-everywhere tiling of $\Omega$ invariant
under the same subgroup, viewed as a full-rank lattice in $\R^3$.
This is also excluded by \cref{geom:no-periodic}, proving both
claims.
\end{proof}

\begin{remark}[Finite selection of the tile data]
The finite data in the construction can be selected by terminating
finite searches. The primes can be selected successively by trial division. The tiling
tests involve maps between finite sets. Lemmas~\ref{stack:partition}
and~\ref{geom:random-representatives} give explicit inequalities
under which a suitable finite coloring or residue transversal exists;
exhaustive search then selects one. This describes
the finite tile itself. The infinite tiling that proves existence is
instead specified by the last-nonzero-digit function. No practical bound
on the size of the finite construction is required.
\end{remark}
